\section*{Bab \ref{ch:b1:erosion-transport} --- Erosi dan transpor}

\begin{solution}{exo:b1:erosion-transport:1}
\qty{0.4}{g.L^{-1}} adalah \qty{0.4}{kg.m^{-3}}, sehingga fluksnya $500 \times 0.4 =
\qty{200}{kg.s^{-1}}$. Setahun dengan \num{365.25}~hari berlangsung
\qty{3.156e7}{s}, sehingga fluks tahunannya adalah \qty{6.31e9}{kg}, atau
\qty{6.31}{Mt}.
\end{solution}

\begin{solution}{exo:b1:erosion-transport:2}
$w_s = R g d^2/(18\nu) = 1.6548 \times 9.81 \times (\num{30e-6})^2/(18 \times
\num{1.0034e-6}) = \qty{8.09e-4}{m.s^{-1}}$, sekitar \qty{0.8}{mm} per detik. Bilangan
Reynolds butirannya adalah $w_s d/\nu = \num{8.09e-4} \times
\num{30e-6}/\num{1.0034e-6} = 0.024$, jauh di bawah 1: hukum Stokes berlaku. Ia
tenggelam melalui \qty{2}{m} dalam $2/\num{8.09e-4} = \qty{2472}{s}$, sekitar
\qty{41}{min}.
\end{solution}

\begin{solution}{exo:b1:erosion-transport:3}
$\tau_b = \rho g h S = 998.2 \times 9.81 \times 0.5 \times \num{4e-4} =
\qty{1.96}{Pa}$. Dengan $(\rho_s - \rho) g = 1651.8 \times 9.81 =
\qty{16204}{N.m^{-3}}$: untuk \qty{1}{mm}, $\theta = 1.96/(16204 \times 0.001) =
0.121$, hampir empat kali ambang \num{0.0314}, sehingga pasir bergerak; untuk
\qty{10}{mm}, $\theta = 0.0121$, seperlima dari ambang \num{0.0556}, sehingga
kerikil tetap diam. Aliran membawa pasirnya di atas dasar kerikil yang tidak bergerak.
\end{solution}

\begin{solution}{exo:b1:erosion-transport:4}
(a)~$12.6/3.9 = 3.2$: sungai membawa sekitar tiga kali lebih banyak dalam suspensi
daripada dalam larutan. (b)~Waduk besar: $0.20 \times 14 = \qty{2.8}{Gt}$ setahun;
semuanya: $0.26 \times 14 = \qty{3.64}{Gt}$ setahun. (c)~Fluks terlarut adalah
\qty{3.9e15}{g} setahun pada \qty{120}{g.m^{-3}}, sehingga air yang membawanya adalah
$\num{3.9e15}/120 = \qty{3.25e13}{m^3}$, sekitar \qty{32500}{km^3} setahun. \omterm{def:b1:erosion-transport:load}{Muatan
melayang} dalam air itu adalah $\qty{1.26e16}{g}/\qty{3.25e16}{L} =
\qty{0.39}{g.L^{-1}}$.
\end{solution}

\begin{solution}{exo:b1:erosion-transport:5}
(a)~Penampang memiliki luas $10 \times 2 = \qty{20}{m^2}$ dan keliling basah
$10 + 2 \times 2 = \qty{14}{m}$, sehingga $R_h = 20/14 = \qty{1.43}{m}$.
(b)~$\tau_b = 998.2 \times 9.81 \times 1.43 \times 0.001 = \qty{14.0}{Pa}$,
dibanding $\qty{19.6}{Pa}$ untuk saluran yang sangat lebar berkedalaman \qty{2}{m}: rasionya
$R_h/h = 0.714$, tebing mengambil sisa gesekan. (c)~Untuk saluran persegi dengan lebar $b$,
$R_h = bh/(b + 2h) = h/(1 + 2h/b)$, yang mendekati $h$ ketika $b$ bertambah.
\end{solution}

\begin{solution}{exo:b1:erosion-transport:6}
(a)~$d_1 = \left(18\nu^2/(Rg)\right)^{1/3} = (18/(Rg))^{1/3}\,\nu^{2/3}$, dan $R$
dan $g$ tidak bergantung pada viskositas. (b)~$\nu = \num{1.3059e-3}/999.7 =
\qty{1.306e-6}{m^2.s^{-1}}$ pada \qty{10}{\celsius}, dibanding
\qty{1.0034e-6}{m^2.s^{-1}} pada \qty{20}{\celsius}: rasionya \num{1.302},
dan $1.302^{2/3} = 1.192$, sehingga $d_1 = 0.104 \times 1.192 = \qty{0.124}{mm}$.
(c)~Dalam rentang Stokes $w_s$ berbanding terbalik dengan $\nu$: butiran tenggelam
$1.302$ kali lebih lambat, sekitar \qty{30}{\percent}. Air dingin menahan sedimen halus
dalam suspensi lebih lama: sungai di musim dingin membawa lanaunya sedikit lebih jauh
untuk aliran yang sama.
\end{solution}

\begin{solution}{exo:b1:erosion-transport:7}
Dengan $R g/(18\nu) = \qty{8.99e5}{m^{-1}.s^{-1}}$ dan hukum gabungan:
\begin{center}
\begin{tabular}{@{}lccc@{}}
\hline
$d$ & Stokes (\si{m.s^{-1}}) & gabungan (\si{m.s^{-1}}) & rasio\\
\hline
\qty{0.1}{mm} & \num{8.99e-3} & \num{7.53e-3} & \num{1.19}\\
\qty{0.5}{mm} & \num{0.225} & \num{0.0711} & \num{3.16}\\
\qty{2}{mm} & \num{3.60} & \num{0.197} & \num{18.3}\\
\hline
\end{tabular}
\end{center}
Rasio tumbuh seiring ukuran, karena bilangan Reynolds butiran tumbuh: hukum Stokes,
yang hanya berlaku untuk bilangan Reynolds kecil, meningkat sebagai $d^2$, sementara
seret jejak memperlambat butiran kasar ke kecepatan yang hanya tumbuh sebagai $\sqrt d$.
\end{solution}

\begin{solution}{exo:b1:erosion-transport:8}
(a)~$(Rg/\nu^2)^{1/3} = \qty{25263}{m^{-1}}$. Pada \qty{0.2}{mm}, $D_* = 5.05$
dan $\theta_c = 0.3/7.06 + 0.055(1 - e^{-0.101}) = 0.0425 + 0.0053 =
0.0478$; $\tau_c = 0.0478 \times 16204 \times \num{2e-4} = \qty{0.155}{Pa}$.
Pada \qty{20}{mm}, $D_* = 505$ dan $\theta_c = 0.3/607 + 0.055(1 - e^{-10.1}) =
0.0005 + 0.0550 = 0.0555$; $\tau_c = 0.0555 \times 16204 \times 0.02 =
\qty{18.0}{Pa}$. (b)~$h = \tau_c/(\rho g S)$ dengan $\rho g S =
\qty{19.58}{Pa.m^{-1}}$: \qty{7.9}{mm} air menggerakkan pasir halus, \qty{0.92}{m}
kerikil.
\end{solution}

\begin{solution}{exo:b1:erosion-transport:9}
Cekungan kehilangan \qty{0.8}{Mt} = \qty{8e5}{t} setahun, atau $\num{8e5}/2000 =
\qty{400}{t.km^{-2}.yr^{-1}}$. Dengan \cref{prop:b1:erosion-transport:denudation},
$D = \qty{8e8}{kg}/(2650 \times \qty{2e9}{m^2}) = \qty{1.51e-4}{m}$ setahun, atau
\qty{151}{m/Ma}. Itu sekitar tiga kali median cekungan yang diukur dengan berilium-10,
\qty{54}{m/Ma}, meskipun di bawah rata-rata mereka \qty{218}{m/Ma}: cekungan yang
cukup cepat, seperti di negeri yang curam atau aktif secara tektonik.
\end{solution}

\begin{solution}{exo:b1:erosion-transport:10}
(a)~$u_* = \sqrt{g h S} = \sqrt{9.81 \times 3 \times \num{5e-4}} =
\qty{0.1213}{m.s^{-1}}$ dan $\tau_b = \rho u_*^2 = \qty{14.7}{Pa}$.
(b)~Suspensi membutuhkan $w_s < u_* = \qty{0.1213}{m.s^{-1}}$. Hukum gabungan
memberi \qty{0.126}{m.s^{-1}} pada \qty{1}{mm}; menyelesaikan
$w_s(d) = 0.1213$ dengan coba-coba memberi $d = \qty{0.94}{mm}$. (c)~Untuk kerikil
$\theta_c \approx 0.0555$, sehingga $\tau_c = \qty{14.7}{Pa}$ untuk $d \approx
14.7/(0.0555 \times 16204) = \qty{16.3}{mm}$; pada ukuran itu kurva ambang memang
memberi $\theta_c = 0.0555$. (d)~Butiran yang lebih kasar dari sekitar \qty{16}{mm}
tetap di dasar dan membentuk lantai kerikil yang stabil; kerikil yang lebih halus dan
pasir yang lebih kasar dari sekitar satu milimeter berpindah sebagai \omterm{def:b1:erosion-transport:load}{muatan dasar};
pasir yang lebih halus, lanau dan lempung berpindah dalam suspensi.
\end{solution}

\begin{solution}{exo:b1:erosion-transport:11}
(a)~$0.30 \times \qty{1000}{m} = \qty{300}{m}$ batuan, dirata-ratakan di seluruh
wilayah. (b)~Pemantulannya adalah $300 \times 2700/3300 = \qty{245}{m}$. (c)~Permukaan
rerata hanya turun sebesar $300 - 245 = \qty{55}{m}$. (d)~Puncak-puncak, yang tidak
tererosi, terangkat bersama seluruh wilayah: kini mereka berdiri \qty{245}{m} lebih
tinggi dari sebelumnya. Pemotongan lembah dapat menaikkan puncak-puncak sebuah pegunungan.
\end{solution}

\begin{solution}{exo:b1:erosion-transport:12}
(a)~$0.9 \times \qty{5e8}{kg}/\qty{1300}{kg.m^{-3}} = \qty{346000}{m^3}$ setahun,
atau \qty{0.69}{\percent} dari waduk. (b)~Setengah waduk, \qty{25e6}{m^3}, terisi dalam
$\num{25e6}/\num{346000} = \num{72}\text{ tahun}$ dengan laju ini. (c)~Ketika waduk
terisi, volumenya, dan dengan demikian waktu tinggal air, berkurang, dan butiran yang
lebih halus lolos: efisiensi penangkapannya turun. Sementara itu endapan yang lebih
tua terkubur dan terpadatkan, mengeluarkan air, sehingga massa jenis curahnya naik dan
setiap ton memakan ruang yang lebih kecil. Kedua efek ini memperlambat kehilangan
volume ketika waduk menua.
\end{solution}

\begin{solution}{pb:b1:erosion-transport:1}
\textbf{1.} Setahun air adalah $1.2 \times \num{3.156e7} = \qty{3.79e7}{m^3}$;
tersebar di atas \qty{8e7}{m^2}, sebuah lapisan setebal \qty{0.473}{m}: limpasan
sebesar \qty{473}{mm} setahun.

\textbf{2.} $1.2 \times \qty{0.150}{kg.m^{-3}} = \qty{0.18}{kg.s^{-1}}$, atau
\qty{5680}{t} setahun.

\textbf{3.} $1.2 \times 0.080 = \qty{0.096}{kg.s^{-1}}$, atau \qty{3030}{t} setahun.

\textbf{4.} \omterm{def:b1:erosion-transport:load}{Muatan dasar} adalah \qty{568}{t} setahun, dan totalnya $5680 + 3030 +
568 = \qty{9278}{t}$ setahun.

\textbf{5.} Konsentrasi sedimen melayang tinggi selama beberapa hari banjir, ketika
debit juga tinggi, dan rendah sepanjang sisa waktu. Rerata waktu memberi bobot yang
sama pada setiap hari, sehingga ia didominasi oleh banyaknya hari aliran rendah dan
konsentrasi rendah. Menurut \cref{prop:b1:erosion-transport:flux}, massa tahunan
adalah $\int Q\,C\,\mathrm dt$; rerata berbobot aliran adalah integral ini dibagi
$\int Q\,\mathrm dt$, sehingga rerata berbobot aliranlah yang harus dikalikan dengan
debit rerata.

\textbf{6.} $3030/9278 = \num{0.33}$, sepertiga. Untuk sungai-sungai dunia,
$3.9/(12.6 + 3.9) = \num{0.24}$, dan lebih kecil jika \omterm{def:b1:erosion-transport:load}{muatan dasar} dihitung.
Bagian cekungan ini lebih tinggi: perbukitan rendah, tempat erosi mekanis yang lambat
memberi batuan waktu untuk larut, dan batu gamping serta kapur tulis, yang larut sepenuhnya.

\textbf{7.} $9278/80 = \qty{116}{t.km^{-2}.yr^{-1}}$.

\textbf{8.} $D = \qty{9.278e6}{kg}/(2700 \times \qty{8e7}{m^2}) =
\qty{4.30e-5}{m}$ setahun, \qty{43.0}{m/Ma}.

\textbf{9.} $43.0/54 = 0.80$: sedikit di bawah median cekungan yang diukur, dan lima
kali di bawah rata-rata mereka. Cekungan yang agak lambat, seperti yang diharapkan
untuk perbukitan rendah di wilayah tenang.

\textbf{10.} Pengukuran debit merata-ratakan beberapa tahun pengukuran, dan bergantung
pada banjir yang terjadi di dalamnya. Berilium-10 merata-ratakan waktu yang diperlukan
untuk mengikis beberapa meter tempat ia dihasilkan: pada \qty{43}{m/Ma}, setiap meter
membutuhkan \num{23000} tahun, sehingga rerata mencakup puluhan ribu tahun. Keduanya
dapat berbeda karena banjir besar yang jarang, yang tidak ada dalam rekaman pendek,
karena pertanian dan pembukaan lahan telah meningkatkan erosi dalam beberapa abad
terakhir, dan karena sedimen tersimpan di dataran banjir dan diambil kembali darinya.

\textbf{11.} Sebagian \omterm{def:b1:erosion-transport:load}{muatan terlarut} tidak berasal dari batuan: banyak bikarbonat
membawa karbon yang diambil dari udara atau tanah, dan hujan membawa garam laut.
Menghitungnya akan memberi cekungan kredit atas batuan yang belum hilang.

\textbf{12.} $\num{1.5e6}/1.2 = \qty{1.25e6}{s}$, \num{14.5} hari.

\textbf{13.} Hukum gabungan memberi \qty{7.53e-3}{m.s^{-1}}: butiran itu jatuh
\qty{8}{m} dalam \qty{1062}{s}, sekitar \qty{18}{min}. Ia terperangkap.

\textbf{14.} $w_s = \num{8.99e5} \times (\num{2e-5})^2 =
\qty{3.60e-4}{m.s^{-1}}$: \qty{8}{m} dalam \qty{22200}{s}, sekitar \num{6.2}~jam.
Terperangkap.

\textbf{15.} $w_s = \num{8.99e5} \times (\num{1e-6})^2 =
\qty{8.99e-7}{m.s^{-1}}$: \qty{8}{m} membutuhkan \qty{8.9e6}{s}, \num{103}~hari, tujuh
kali waktu tinggal. Sebagian besar darinya melewati bendungan bersama air.

\textbf{16.} Butiran harus jatuh \qty{8}{m} dalam \qty{1.25e6}{s}, pada kecepatan
paling sedikit $\qty{6.4e-6}{m.s^{-1}}$. Menurut hukum Stokes $d = \sqrt{18\nu
w_s/(Rg)} = \sqrt{\num{6.4e-6}/\num{8.99e5}} = \qty{2.7}{\micro m}$: lempung halus.
Butiran yang lebih kasar terperangkap di mana pun mereka masuk.

\textbf{17.} $568 + 0.7 \times 5680 = \qty{4544}{t}$ setahun, yang mengisi
$\num{4.544e6}/1300 = \qty{3496}{m^3}$ setahun, \qty{0.233}{\percent} dari waduk.
Setengah dari waduk itu, \qty{7.5e5}{m^3}, terisi dalam sekitar \num{215} tahun.

\textbf{18.} $100/43.0 = \qty{2.33}{Ma}$.

\textbf{19.} $1 - 2700/3300 = 0.18$: permukaan turun hanya sebesar
\qty{18}{\percent} dari ketebalan yang dihilangkan.

\textbf{20.} $100/0.1818 = \qty{550}{m}$ batuan.

\textbf{21.} Seluruh wilayah naik sebesar $550 \times 2700/3300 = \qty{450}{m}$
akibat pemantulan. Puncak-puncak yang nyaris tidak tererosi terangkat bersamanya, hingga
\qty{450}{m}, sementara lembah-lembah makin dalam: relief tumbuh sebelum ia menurun.

\textbf{22.} \omterm{def:b1:erosion-transport:denudation}{Laju denudasi} bergantung terutama pada lereng. Ketika perbukitan diturunkan
dan lerengnya melandai, laju akan turun, kecuali pengangkatan atau perubahan iklim
memperkuatnya kembali: waktu yang dihitung dengan laju hari ini adalah batas bawah.

\textbf{23.} $550/42.95 = 12.80$. Pada laju hari ini dan dengan pemantulan isostatik,
menurunkan perbukitan Varden sebesar \qty{100}{m} membutuhkan \textbf{12.8}~juta tahun.
\end{solution}
